\documentclass[11pt]{article}

\usepackage{amsmath, amssymb}
\usepackage{geometry}
\usepackage{enumitem}
\usepackage{hyperref}

\geometry{margin=1in}

\title{\textbf{CSE400 — Lecture 11}\\
Transformation of Random Variables\\
\large Exam-Oriented Lecture Scribe for Revision}
\author{
Course: CSE400 — Fundamentals of Probability in Computing\\
Instructor: Dhaval Patel, Associate Professor, SEAS, Ahmedabad University\\
Lecture Date: February 10, 2026
}
\date{}

\begin{document}

\maketitle

This scribe reconstructs the lecture content strictly from the lecture slides and preserves the structure, results, and solution steps presented.

\hrule
\vspace{0.5em}

\section*{1. Lecture Outline}

The lecture covers:

\begin{enumerate}
    \item \textbf{Transformation of Random Variables}
    \begin{itemize}
        \item Learning techniques to transform random variables.
    \end{itemize}

    \item \textbf{Function of Two Random Variables}
    \begin{itemize}
        \item Joint transformations and derived distributions.
    \end{itemize}

    \item \textbf{Illustrative Example}
    \begin{itemize}
        \item Detailed derivation for the case:
        \[
        Z = X + Y
        \]
    \end{itemize}
\end{enumerate}

\hrule
\vspace{0.5em}

\section*{2. Transformation of Random Variables}

The lecture considers transformation of a random variable where:

\begin{itemize}
    \item A random variable \(X\) is known.
    \item A new random variable \(Y\) is defined as a function of \(X\):
    \[
    Y = g(X)
    \]
\end{itemize}

The goal is to obtain the distribution of \(Y\).

\subsection*{Method Used: CDF-Based Transformation}

The cumulative distribution function (CDF) of \(Y\) is derived using probability relations.

\subsubsection*{Step-by-step derivation}

\begin{enumerate}
    \item Define CDF of \(Y\):
    \[
    F_Y(y) = P(Y \le y)
    \]

    \item Substitute transformation:
    \[
    = P(g(X) \le y)
    \]

    \item Express event in terms of \(X\):
    \[
    = P(X \le g^{-1}(y))
    \]
    (depending on monotonicity of the transformation).

    \item Use the CDF of \(X\):
    \[
    = F_X(g^{-1}(y))
    \]

    \item Obtain PDF of \(Y\) by differentiation:
    \[
    f_Y(y) = \frac{d}{dy} F_Y(y)
    \]
\end{enumerate}

This procedure uses the known PDF or CDF of \(X\) to derive that of \(Y\).

\hrule
\vspace{0.5em}

\section*{3. Example: Transformation of a Uniform Random Variable}

An example involves a uniformly distributed random variable.

From the lecture slides:

\begin{itemize}
    \item A random variable is uniformly distributed over an interval.
    \item A transformation is applied:
    \[
    Y = \sin(\cdot)
    \]
    (exact argument as given in slides).
\end{itemize}

Steps followed:

\begin{enumerate}
    \item Start with known PDF of the original variable.
    \item Apply transformation relation.
    \item Express transformed variable in terms of original variable.
    \item Use probability relation:
    \[
    F_Y(y) = P(Y \le y)
    \]
    \item Convert inequality into original variable’s domain.
    \item Use known distribution of original variable.
    \item Differentiate to obtain PDF.
\end{enumerate}

The lecture maintains symbolic manipulation from original to transformed distribution.

\hrule
\vspace{0.5em}

\section*{4. Function of Two Random Variables}

The lecture then considers transformations involving two random variables.

\begin{itemize}
    \item Two random variables:
    \[
    X, Y
    \]
    \item Possibly independent.
    \item A new variable defined as:
    \[
    Z = h(X, Y)
    \]
\end{itemize}

The distribution of derived variables is determined using the joint PDF.

\hrule
\vspace{0.5em}

\section*{5. Worked Example: Distribution of \(Z = X + Y\)}

\subsection*{Problem Setup}

A new random variable is defined as:
\[
Z = X + Y
\]

Goal: find the distribution of \(Z\).

\subsection*{Step 1: Express CDF}

\[
F_Z(z) = P(Z \le z)
\]

Substituting:
\[
= P(X + Y \le z)
\]

\subsection*{Step 2: Joint Probability Region}

The event corresponds to all pairs \((x, y)\) satisfying:
\[
x + y \le z
\]

Thus,
\[
F_Z(z) = P((X, Y) \text{ lie in this region})
\]

\subsection*{Step 3: Joint PDF Integration}

The CDF becomes:
\[
F_Z(z)
= \iint_{x + y \le z} f_{X,Y}(x,y)\, dy\, dx
\]

Integration occurs over the region defined by:
\[
x + y \le z
\]

\subsection*{Step 4: Ordered Integral Form}

Slides express:
\[
F_Z(z)
= \int \int f_{X,Y}(x,y)\, dy\, dx
\]
with limits determined by the region.

\subsection*{Step 5: Final Objective}

Once CDF is obtained,
\[
f_Z(z) = \frac{d}{dz}F_Z(z)
\]
gives the PDF of \(Z\).

\hrule
\vspace{0.5em}

\section*{6. Logical Flow of Lecture}

Lecture sequence:

\begin{enumerate}
    \item Review transformation of a single random variable.
    \item Apply CDF-based method.
    \item Demonstrate with an example.
    \item Extend to two variables.
    \item Derive distribution for \(Z = X + Y\).
    \item Express solution via joint PDF integration.
\end{enumerate}

\hrule
\vspace{0.5em}

\section*{7. Key Results for Exams}

Important results:

\begin{itemize}
    \item Transformation uses:
    \[
    F_Y(y) = P(Y \le y)
    \]
    \item Convert event into original variables.
    \item Use known distributions.
    \item Differentiate to obtain PDFs.
    \item For two variables:
    \[
    F_Z(z) = P(X + Y \le z)
    \]
    \item Use joint PDF integration over region.
    \item PDF obtained by differentiating the CDF.
\end{itemize}

\hrule
\vspace{0.5em}

\section*{End of Lecture 11 Scribe}

This document contains all derivations and solution steps presented in lecture slides and is structured for direct exam revision.

\end{document}
```
